
\documentclass{article}
\usepackage{colortbl}
\usepackage{makecell}
\usepackage{multirow}
\usepackage{supertabular}

\begin{document}

\newcounter{utterance}

\centering \large Interaction Transcript for game `wordle\_withcritic\_noexp', experiment `medium\_frequency\_words\_clue\_with\_critic', episode 3 with reflexion{-}llama{-}3.1{-}8b{-}instruct.
\vspace{24pt}

{ \footnotesize  \setcounter{utterance}{1}
\setlength{\tabcolsep}{0pt}
\begin{supertabular}{c@{$\;$}|p{.15\linewidth}@{}p{.15\linewidth}p{.15\linewidth}p{.15\linewidth}p{.15\linewidth}p{.15\linewidth}}
    \# & $\;$A & \multicolumn{4}{c}{Game Master} & $\;\:$B\\
    \hline

    \theutterance \stepcounter{utterance}  
    & & \multicolumn{4}{p{0.6\linewidth}}{
        \cellcolor[rgb]{0.9,0.9,0.9}{
            \makecell[{{p{\linewidth}}}]{
                \texttt{\tiny{[P1$\langle$GM]}}
                \texttt{You have attempted this same task before and did not succeed. These are the plans you wrote after those attempts.} \\
\texttt{Trial \#1: New plan: The last episode got aborted because "sweetly" is not a five{-}letter word. I should have paid closer attention to the length of the word. Next time, I will make sure to produce a five{-}letter word that does not contain red letters (s, w, e, t), for example, "fable".} \\
\\ 
\texttt{Trial \#2: New plan: The last episode got aborted because "fable" is not a valid guess based on the clue "sweetheart". I should have paid closer attention to the clue and the length of the word. The clue suggests a word related to love or affection, and "fable" does not fit this context. Next time, I will make sure to produce a five{-}letter word that is related to the clue and does not contain red letters (s, w, e, t), for example, "honey".} \\
\\ 
\\ 
\texttt{You are a language wizard who likes to guess words by using the given rules.} \\
\\ 
\texttt{Welcome to Wordle! You have six attempts to guess the target word, a valid English word of five lowercase letters (a{-}z). Please use the tag "guess:" to provide your guess.} \\
\\ 
\texttt{To help you make an informed guess, you will receive a clue for the word, such as} \\
\texttt{clue: "snowy white"} \\
\\ 
\texttt{Here is an example guess based on the clue: } \\
\texttt{guess: apple} \\
\\ 
\texttt{I will then indicate whether I agree or disagree with your guess, but agreeing with a guess does not confirm its correctness. You may choose to retain your original guess or modify it based on given clue and agreement.} \\
\\ 
\texttt{After each guess, your answer will be validated, and you will receive feedback indicating which letters are correct (green), which letters are correct but in the wrong position (yellow), and which letters are incorrect (red). This feedback can be useful in determining which letters to include or exclude in your next guess.} \\
\\ 
\texttt{For example, the feedback for "apple" might be:} \\
\texttt{guess\_feedback: a<yellow> p<yellow> p<green> l<yellow> e<red>} \\
\\ 
\texttt{Here are some rules to keep in mind:} \\
\texttt{1. Please do not generate any guess feedback or clue on your own} \\
\texttt{2. Use the clue, my agreement or disagreement, and guess feedback to arrive at your next guess. } \\
\\ 
\texttt{Do not generate any other text.} \\
\\ 
\texttt{Let's begin with your first guess.} \\
\\ 
\texttt{clue: sweetheart} \\
            }
        }
    }
    & \\ \\

    \theutterance \stepcounter{utterance}  
    & \multicolumn{4}{p{0.6\linewidth}}{
        \cellcolor[rgb]{0.9,0.9,0.9}{
            \makecell[{{p{\linewidth}}}]{
                \texttt{\tiny{[P1$\rangle$GM]}}
                \texttt{guess: honey} \\
            }
        }
    }
    & & \\ \\

    \theutterance \stepcounter{utterance}  
    & & \multicolumn{4}{p{0.6\linewidth}}{
        \cellcolor[rgb]{0.9,0.9,0.9}{
            \makecell[{{p{\linewidth}}}]{
                \texttt{\tiny{[GM$\rangle$P2]}}
                \texttt{You have attempted this same task before and did not succeed. These are the plans you wrote after those attempts.} \\
\texttt{Trial \#1: New plan: I got stuck in a loop of responding with "no" without considering the guess feedback. In my last attempt, I didn't use the guess feedback to inform my decision. Next time, I will carefully examine the guess feedback and use it to decide whether the given guess matches the clue. If the guess feedback contains green or yellow letters, I will take that into account when making my decision. For example, if the guess feedback contains a green letter, I will check if that letter is in the correct position in the guess. If the guess feedback contains a yellow letter, I will check if that letter is present in the clue but not in the correct position in the guess.} \\
\\ 
\texttt{Trial \#2: New plan: In my last attempt, I got stuck in a loop of responding with "no" without considering the guess feedback. However, I also failed to follow the game rules by not sticking to the given format while responding. Next time, I will make sure to carefully examine the guess feedback and use it to decide whether the given guess matches the clue. If the guess feedback contains green or yellow letters, I will take that into account when making my decision. Additionally, I will ensure that my response is in lowercase letters and follows the specified format, which is "agreement: yes or no".} \\
\\ 
\\ 
\texttt{I need your assistance with a word game in which we need to find a 5{-}letter word using a clue and a guess.} \\
\\ 
\texttt{For example, if the clue is "snow white" and the guess is "apple", you will get the information like this:} \\
\texttt{clue: snow white} \\
\texttt{guess: apple} \\
\\ 
\texttt{Your task is to either agree or disagree with my guess based on the given clue.} \\
\\ 
\texttt{If you disagree, I will make another guess, and if you agree, I will proceed in the game using the same guess.} \\
\\ 
\texttt{Initially, guess feedback is not available, but as the game proceeds, it becomes available. The guess feedback is not for current guess, but rather an indication of what letters can be present in the current guess. A letter highlighted in green indicates that it is a correct letter in the correct position, while yellow indicates that it is a correct letter in the wrong position, and red indicates a wrong letter. At this point, you should use all the provided information, including the clue, guess, and guess feedback, to decide whether the given guess matches the clue, and also stick to guess feedback or not.} \\
\\ 
\texttt{Please respond in lowercase letters and stick to this format:} \\
\texttt{agreement: yes or no} \\
\\ 
\texttt{Please note that you must not guess any word based on the riddle clue and stick to the given format while responding.} \\
\\ 
\texttt{Do not generate any other text.} \\
\\ 
\texttt{Let's begin!} \\
\\ 
\texttt{clue: sweetheart} \\
\texttt{guess: honey} \\
            }
        }
    }
    & \\ \\

    \theutterance \stepcounter{utterance}  
    & & & \multicolumn{4}{p{0.6\linewidth}}{
        \cellcolor[rgb]{0.9,0.9,0.9}{
            \makecell[{{p{\linewidth}}}]{
                \texttt{\tiny{[GM$\langle$P2]}}
                \texttt{agreement: no} \\
            }
        }
    }
     \\ \\

    \theutterance \stepcounter{utterance}  
    & & \multicolumn{4}{p{0.6\linewidth}}{
        \cellcolor[rgb]{0.9,0.9,0.9}{
            \makecell[{{p{\linewidth}}}]{
                \texttt{\tiny{[P1$\langle$GM]}}
                \texttt{clue: sweetheart} \\
\texttt{guess\_agreement: no} \\
\texttt{Provide your response only in this format.} \\
\texttt{guess: word} \\
            }
        }
    }
    & \\ \\

    \theutterance \stepcounter{utterance}  
    & \multicolumn{4}{p{0.6\linewidth}}{
        \cellcolor[rgb]{0.9,0.9,0.9}{
            \makecell[{{p{\linewidth}}}]{
                \texttt{\tiny{[P1$\rangle$GM]}}
                \texttt{guess: honey} \\
            }
        }
    }
    & & \\ \\

    \theutterance \stepcounter{utterance}  
    & & & \multicolumn{2}{p{0.3\linewidth}}{
        \cellcolor[rgb]{0.9,0.9,0.9}{
            \makecell[{{p{\linewidth}}}]{
                \texttt{\tiny{[GM$|$GM]}}
                \texttt{attempts = 1} \\
\texttt{target\_word = flame} \\
\texttt{target\_word\_clue = sweetheart} \\
\texttt{guess = honey} \\
\texttt{guess\_feedback = h<red> o<red> n<red> e<yellow> y<red>} \\
            }
        }
    }
    & & \\ \\

    \theutterance \stepcounter{utterance}  
    & & \multicolumn{4}{p{0.6\linewidth}}{
        \cellcolor[rgb]{0.9,0.9,0.9}{
            \makecell[{{p{\linewidth}}}]{
                \texttt{\tiny{[P1$\langle$GM]}}
                \texttt{guess\_feedback: h<red> o<red> n<red> e<yellow> y<red>} \\
\\ 
\texttt{Provide your response only in this format.} \\
\texttt{guess: word} \\
            }
        }
    }
    & \\ \\

    \theutterance \stepcounter{utterance}  
    & \multicolumn{4}{p{0.6\linewidth}}{
        \cellcolor[rgb]{0.9,0.9,0.9}{
            \makecell[{{p{\linewidth}}}]{
                \texttt{\tiny{[P1$\rangle$GM]}}
                \texttt{guess: lovely} \\
            }
        }
    }
    & & \\ \\

    \theutterance \stepcounter{utterance}  
    & & & \multicolumn{2}{p{0.3\linewidth}}{
        \cellcolor[rgb]{0.9,0.9,0.9}{
            \makecell[{{p{\linewidth}}}]{
                \texttt{\tiny{[GM$|$GM]}}
                \texttt{Error: The length of the guessed word is not 5.} \\
            }
        }
    }
    & & \\ \\

    \theutterance \stepcounter{utterance}  
    & & & \multicolumn{2}{p{0.3\linewidth}}{
        \cellcolor[rgb]{0.9,0.9,0.9}{
            \makecell[{{p{\linewidth}}}]{
                \texttt{\tiny{[GM$|$GM]}}
                \texttt{game\_result = ABORT} \\
            }
        }
    }
    & & \\ \\

    \theutterance \stepcounter{utterance}  
    & \multicolumn{4}{p{0.6\linewidth}}{
        \cellcolor[rgb]{0.95,0.95,0.95}{
            \makecell[{{p{\linewidth}}}]{
                \texttt{\tiny{[P2$\rightarrow$P1]}}
                \texttt{Post{-}episode Reflexion (not sent to GM):} \\
\texttt{New plan: The last episode got aborted because "lovely" is not a valid guess based on the clue "sweetheart" and the previous guess "honey". I should have paid closer attention to the clue and the previous feedback. The clue suggests a word related to love or affection, and "honey" already contains the letter "h" and "e" in the correct position. The previous feedback also indicates that "e" is in the correct position. Next time, I will make sure to produce a five{-}letter word that is related to the clue, contains the letter "h" and "e" in the correct position, and does not contain red letters (s, w, t), for example, "hearts".} \\
            }
        }
    }
    & & \\ \\

    \theutterance \stepcounter{utterance}  
    & & & \multicolumn{4}{p{0.6\linewidth}}{
        \cellcolor[rgb]{0.95,0.95,0.95}{
            \makecell[{{p{\linewidth}}}]{
                \texttt{\tiny{[P1$\rightarrow$P2]}}
                \texttt{Post{-}episode Reflexion (not sent to GM):} \\
\texttt{New plan: In my last attempt, I got stuck in a loop of responding with "no" without considering the guess feedback, and I also failed to follow the game rules by not sticking to the given format while responding. However, I also failed to use the clue information in my decision. Next time, I will make sure to carefully examine the guess feedback and use it to decide whether the given guess matches the clue. If the guess feedback contains green or yellow letters, I will take that into account when making my decision. Additionally, I will ensure that my response is in lowercase letters and follows the specified format, which is "agreement: yes or no". Furthermore, I will also use the clue information to inform my decision, considering the context and meaning of the clue. For example, if the clue is "snow white" and the guess is "apple", I will consider the fact that "apple" is not related to "snow white" and the guess feedback will also support this decision.} \\
            }
        }
    }
     \\ \\

\end{supertabular}
}

\end{document}
